Cryptographic protocols such as Transport Layer Security (TLS), Secure Shell (SSH), and Virtual Private Networks (VPNs) rely on asymmetric cryptography to securely establish shared secrets between previously untrusted parties. These cryptographic systems are based on the assumption that specific mathematical problems, notably integer factorization and discrete logarithms, are too hard for classical computers to solve efficiently.

However, this assumption may not hold in the long term. Rapid advances in quantum computing research suggest that large-scale, fault-tolerant quantum computers may become viable within the coming decades \cite{Swayne2025QuantumRoadmaps}. These are computers capable of leveraging the laws of quantum mechanics to perform quantum algorithms that achieves significant computational speedups for specific classes of problems in comparison to their classical counterparts. Such machines would be able to fundamentally break the asymmetric cryptography primitives used today. 

Two quantum algorithms stand out in this cryptographic context. Peter Shor's algorithm can efficiently factor large integers and compute discrete logarithms \cite{shors}, and Lov Grover's algorithm \cite{grovers}, which provides a quadratic speedup to brute-force key searches. The former directly breaks the public-key primitives underlying RSA, Diffie-Hellman and elliptic-curve cryptography, while the latter reduces the effective security level of symmetric ciphers and hash functions by half.

Once quantum computers reach this critical scale, they will likely be capable of decrypting nearly all historical and contemporary communications secured with current public-key systems. Moreover, adversaries can already capture encrypted traffic today with the intent of decrypting it later when quantum capabilities become available, a strategy commonly referred to as ``harvest now, decrypt later'' \cite{hndl}. This threat makes it essential to begin the transition toward quantum-resistant cryptography as soon as possible, since deploying new cryptographic standards across the Internet ecosystem will take many years.

To stay ahead of this threat, the cryptographic community is actively developing Post-Quantum Cryptography (PQC) algorithms. These are new types of ``quantum-resistant'' mechanisms designed to be secure against both today's computers and tomorrow's quantum machines. The algorithms are based on mathematical problems believed to remain hard for quantum computers, such as those derived from lattices, error-correcting codes, or hash functions. The National Institute of Standards and Technology (NIST) initiated, back in 2016, a global effort to identify and standardize these new methods \cite{nistpqc}, resulting in several promising candidates.

After several years of public evaluation, the NIST Post-Quantum Cryptography Standardization Process selected and standardized ML-KEM \cite{mlkem}, a lattice-based Key-Encapsulation Mechanism (KEM), and is currently in the process of standardizing Hamming Quasi-Cyclic (HQC) \cite{hqc}, a code-based KEM. These schemes are considered primary building blocks for securing Internet communications in the quantum era. 

While the theoretical security of these algorithms has been extensively studied, their practical deployment within real-world communication protocols remains an active research challenge. Post-quantum schemes often have larger key sizes, ciphertexts, and computational overhead compared to classical cryptographic methods. Such differences can significantly affect performance-sensitive protocols that prioritize efficiency.

One such protocol is WireGuard \cite{wireguard}, a modern VPN protocol designed for minimalism, simplicity, and high throughput. WireGuard currently relies on elliptic-curve cryptography, specifically Curve25519, as part of its authenticated key exchange mechanism. While this design offers excellent efficiency on classical hardware, elliptic-curve cryptography is vulnerable to quantum attacks due to its reliance on the discrete logarithm problem.

Integrating PQC into WireGuard therefore presents several unique challenges related to efficiency, compatibility, and code maintainability. These include managing increased message sizes, maintaining low handshake latency, and preserving the protocol's minimalist design philosophy. Existing solutions, such as Rosenpass and PQ-WireGuard prototypes, have demonstrated the feasibility of quantum-resistant VPN connections through external or hybrid key exchanges \cite{rosenpass}\cite{pqwg}\cite{rebar}.

Yet, few studies have systematically compared the performance characteristics of NIST-approved algorithms when integrated into WireGuard's architecture.

This thesis addresses that gap by designing, implementing, and evaluating a hybrid WireGuard handshake that supports two NIST-selected Key-Encapsulation Mechanisms: ML-KEM and Hamming Quasi-Cyclic. The implementation enables an empirical comparison of their latency, computational overhead, memory usage, and message size, and contrasts the results against existing post-quantum WireGuard implementations. By quantifying these trade-offs, the work aims to provide practical insights into the real-world viability of post-quantum VPNs and contribute to the broader effort of securing Internet communications against future quantum threats.

\section{Problem statement and Research Questions}\label{research_questions}
\textit{How does NIST-approved Post-Quantum Key-Encapsulation Mechanisms, specifically ML-KEM and HQC, affect the performance, resource usage, and practical deployability of a post-quantum secure WireGuard handshake compared to existing  approaches?}

\begin{enumerate}[label=\textbf{RQ\arabic*.}]
    \item How do NIST-selected post-quantum key encapsulation mechanisms, specifically ML-KEM and HQC, affect the performance of a WireGuard handshake?

    \item What impact do ML-KEM and HQC have on resource usage in WireGuard, particularly in terms of message size and storage requirements?

    \item To what extent can ML-KEM and HQC be integrated into WireGuard without violating its core design principles, such as low latency and small handshake size?

    \item How do hybrid cryptographic approaches (combining classical Diffie-Hellman with PQC) compare to purely post-quantum approaches in terms of security, performance, and deployability?

    \item How does a hybrid WireGuard implementation using ML-KEM and HQC compare to existing post-quantum solutions in terms of performance and practical deployment?

    \item What trade-offs exist between security level and ciphertext size when selecting PQC algorithms for use in WireGuard?
\end{enumerate}

\section{Scope and Limitations} 
This thesis covers the design, implementation, and empirical evaluation of post-quantum KEM integration into WireGuard's handshake, using ML-KEM and HQC at NIST security levels 3 and 5. The implementation contains both a hybrid mode, which augments the classical handshake with an additional KEM operation, and a pure-KEM mode, which replaces the classical exchange entirely, are implemented and benchmarked. 

Several limitations apply to this work. The implementation is built using \texttt{wireguard-go}, the userspace Go implementation of WireGuard, and does not modify the Linux kernel module. Performance numbers from the kernel implementation would be better due to the absence of TUN interface overhead, so the results here should not be taken as representative of a future kernel-level deployment. 

All benchmarks were conducted on a single hardware platform: two virtual machines running on the same KVM hypervisor. No measurements were taken on physical server hardware, ARM-based processors, or mobile devices.

The benchmark environment uses a virtual network on a single physical host, which eliminates packet loss and real-network jitter. In particular, large handshake messages that exceed the standard MTU are fragmented and reassembled entirely in software, without the reliability problems that IP fragmentation causes on real network paths. Thus, the HQC results in this thesis are measured under more favorable conditions than a real deployment would provide. 

The security analysis of the hybrid and pure-KEM handshake constructions is informal. No formal verification or symbolic proof is provided, and the security arguments in this thesis are based on design reasoning rather than machine-checked proofs. This is consistent with the scope of the related work that this thesis builds on, but it means the security claims are weaker than those of formally analyzed designs.


\section{Thesis Structure} 
The remainder of this thesis is structured as follows.

Chapter 2 provides the necessary background. It begins with classical cryptography, covering classical public-key primitives such as RSA, Diffie-Hellman, and elliptic-curve cryptography, along with key concepts such as key exchange and key encapsulation mechanisms. It then introduces the quantum threat, including a high-level overview of quantum computing and the algorithms that break classical cryptography: Shor's and Grover's algorithms, as well as the ``harvest now, decrypt later'' risk model. Finally, it introduces post-quantum cryptography, covering the main algorithm families with a focus on ML-KEM and HQC, and discusses hybrid cryptographic approaches and the practical challenges of deploying PQC in real protocols. 

Chapter 3 covers the WireGuard protocol. It describes its design principles, the Noise Protocol Framework, the IKpsk2 handshake pattern, the cryptographic primitives in use, the full handshake flow, and the Go implementation. It also discusses the specific challenges that arise when integrating post-quantum algorithms into WireGuard. 

Chapter 4 discusses the related work on post-quantum WireGuard, covering the original PQ-WireGuard by Hülsing et al., the Rosenpass sidecar approach, and the recent Rebar redesign using reinforced KEMs by Hashimoto et al. It also describes concurrent work, covering the ongoing NIST standardization of HQC and the deployment of hybrid ML-KEM in TLS 1.3.

Chapter 5 describes the research method. It covers the literature study, protocol analysis, design decisions, implementation of both the hybrid and pure-KEM handshake variants, measurement methodology, and benchmarking setup. It concludes with a discussion of reliability, validity, and generalizability. 

Chapter 6 presents the results of the benchmark: handshake latency, message sizes, and storage requirements for all nine evaluated variants, compares them against the related work, and describes how they answer the research questions. 

Chapter 7 outlines the contributions of the thesis to literature, practice, and management, examines the MTU constraint in detail, addresses the requirements of full post-quantum authentication, and compares the implementation with PQ-WireGuard, Rosenpass, and Rebar. It also considers the ethical and sustainability implications of the transition to post-quantum cryptography, and aims to direct future research.

Chapter 8 concludes the thesis by summarising the main findings and reflecting on the broader implications for post-quantum VPN deployment.